\section{What TFPT Does \emph{Not} Assume}
%=============================================================================

The following are \emph{not} required:
\begin{itemize}[leftmargin=2em]
    \item[\ding{55}] Slow-roll inflation (specific model)
    \item[\ding{55}] String landscape or anthropic selection
    \item[\ding{55}] Extra dimensions as free parameters
    \item[\ding{55}] Dark matter entirely from axions
    \item[\ding{55}] New particles beyond minimal axion sector
    \item[\ding{55}] Any fitted or adjusted parameters
\end{itemize}

%=============================================================================
\section{Conclusions}
%=============================================================================

TFPT demonstrates that the fine-structure constant $\alpha$ can be derived 
from first principles using eight minimal axioms rooted in topology, geometry, 
and quantum consistency. The key results are:

\begin{enumerate}
    \item The topological coupling $\cthree = 1/(8\pi)$ from 11D Chern--Simons 
        quantization and the geometric scale $\phiz = 1/(6\pi) + 3/(256\pi^4)$ 
        from M\"obius fiber reduction are the only fundamental inputs.
    
    \item The cubic fixed-point equation (CFE) has a unique positive root. With the derived two-defect term $\delta_2=\frac54\delta_{\mathrm{top}}^2$ (holonomy $g=5$), this yields
        $\alpha^{-1}(0)=\num[round-mode=places,round-precision=12]{\TFPTalphaInvGOverFourNum}$ (one-defect truncation: $\num[round-mode=places,round-precision=10]{\TFPTalphaInvSelfNum}$).
    
    \item The same invariants determine the axion--photon coupling 
        $\gagg = -4\cthree$ and predict cosmic birefringence 
        $\beta = \num[round-mode=places,round-precision=4]{\TFPTbetaDegNum}^\circ$, consistent with observations.
    
    \item An E$_8$ cascade organizes all physical scales from Planck to 
        electroweak, with deviations typically below 5\%.
\end{enumerate}

The theory is falsifiable: any significant revision of $\alpha$, $\beta$, 
or the two-loop RG fingerprints would invalidate the framework.

%=============================================================================
\appendix
\section{Numerical Constants (High Precision)}
%=============================================================================

\begin{table}[H]
\centering
\small
\begin{tabular}{@{}ll@{}}
\toprule
\textbf{Quantity} & \textbf{Value}\\
\midrule
$\cthree = 1/(8\pi)$ & \num[round-mode=places,round-precision=12]{\TFPTcThreeNum}\\
$\phitree = 1/(6\pi)$ & \num[round-mode=places,round-precision=12]{\TFPTvarphiTreeNum}\\
$\deltatop = 3/(256\pi^4)$ & \num[round-mode=places,round-precision=12]{\TFPTdeltaTopNum}\\
$\phiz = \phitree + \deltatop$ & \num[round-mode=places,round-precision=12]{\TFPTvarphiZeroNum}\\
$\phiz(\alpha)$ (self-consistent) & \num[round-mode=places,round-precision=12]{0.053170209119023745}\\
$\beta_{\text{rad}} = \phiz/(4\pi)$ & \num[round-mode=places,round-precision=12]{\TFPTbetaRadNum}\\
$\beta_{\text{deg}} = (180/\pi)\,\beta_{\text{rad}}$ & \num[round-mode=places,round-precision=6]{\TFPTbetaDegNum}$^\circ$\\
$g_{a\gamma\gamma}$ & \num[round-mode=places,round-precision=12]{\TFPTgAggNum}\\
$\gamma(0)=g/(g+1)$ (E$_8$) & $\tfrac56$\\
$\lambda$ (E$_8$ exponent) & \num[round-mode=places,round-precision=12]{0.587233190787837868}\\
$M/\Mpl$ ($R^2$ scale) & \num[round-mode=places,round-precision=12]{\TFPTMOverMplNum}\\
$A = 2\cthree^3$ & \num[round-mode=figures,round-precision=8]{1.2598255637968552e-4}\\
$\mathcal{K} = (\bone/2\pi)\ln(1/\phiz)$ & \num[round-mode=places,round-precision=10]{1.91468479457699546}\\
$\alpha^{-1}_{\text{baseline}}$ & \num[round-mode=places,round-precision=9]{137.036501464885822}\\
$\alpha^{-1}_{\text{TFPT}}$ (one-defect) & \num[round-mode=places,round-precision=12]{\TFPTalphaInvSelfNum}\\
$\alpha^{-1}_{\text{TFPT}}$ (two-defect, $g=5$) & \num[round-mode=places,round-precision=12]{\TFPTalphaInvGOverFourNum}\\
$\alpha^{-1}_{\text{TFPT}}$ (two-defect, suite diagnostic) & \num[round-mode=places,round-precision=12]{\TFPTalphaInvTwoDefNum}\\
$\alpha^{-1}_{\text{CODATA 2022}}$ & \num[round-mode=places,round-precision=12]{\TFPTalphaInvCODATANum}\\
Accuracy (ppm; one-defect) & \num[round-mode=places,round-precision=6]{-0.03704}\\
Accuracy (ppm; two-defect, $g=5$) & \num[round-mode=places,round-precision=6]{\TFPTalphaPpmGOverFourNum}\\
Accuracy (ppm; two-defect, suite diagnostic) & \num[round-mode=places,round-precision=6]{\TFPTalphaPpmTwoDefNum}\\
$\xi = \cthree/\phiz$ & \num[round-mode=places,round-precision=9]{0.748303083562547799}\\
\bottomrule
\end{tabular}
\caption{High-precision numerical values of TFPT constants.}
\label{tab:constants}
\end{table}

%=============================================================================
\section{Derivation of the CFE from the Effective Potential}
%=============================================================================

In background-field gauge, the one-loop effective potential for the electromagnetic 
coupling modulus takes the form. Here $U(\alpha)$ is the induced potential for $\sigma$
after rewriting $\alpha=\alpha(\sigma)$:
\begin{equation}
    U(\alpha) = \frac{1}{4}(F_{\mu\nu}F^{\mu\nu})_{\text{eff}} 
        + \text{quantum corrections}
\end{equation}

The full structure, including the Chern--Simons contribution from 11D reduction, is:
\begin{equation}
    U(\alpha) = \frac{A}{4}\alpha^4 - \frac{2}{3}A\cthree^3\alpha^3 
        - 8A\bone\cthree^6 L\,\alpha + \mathcal{O}(\alpha^5)
    \label{eq:potential}
\end{equation}
where $L = \ln(1/\phiz)$ encodes the geometric scale.

The modulus-closure condition $\delta\Gamma_{\mathrm{eff}}/\delta\sigma=0$ reduces to
$\partial U/\partial\alpha = 0$ (monotonic $\alpha(\sigma)$), giving:
\begin{equation}
    \frac{\partial U}{\partial\alpha} = A\alpha^3 - 2A\cthree^3\alpha^2 
        - 8A\bone\cthree^6 L = 0
    \label{eq:stationarity}
\end{equation}

Dividing by $A$ (which is nonzero):
\begin{equation}
    \alpha^3 - 2\cthree^3\alpha^2 - 8\bone\cthree^6 L = 0
    \label{eq:CFE-appendix}
\end{equation}

This is the CFE. The combinatorial factors arise from specific one-loop diagrams:

\textbf{(i) Factor 2 in $2\cthree^3\alpha^2$:}

The cubic term in the effective potential comes from the vacuum polarization 
(``bubble'') diagram with axion insertion:
\begin{center}
\begin{tikzpicture}[scale=0.8]
    \draw[thick, photon] (0,0) -- (1.5,0);
    \draw[thick] (1.5,0) arc (180:0:0.75);
    \draw[thick] (1.5,0) arc (-180:0:0.75);
    \draw[thick, photon] (3,0) -- (4.5,0);
    \draw[thick, dashed] (2.25,0.75) -- (2.25,1.5);
    \node at (2.25,1.8) {\small $a$};
    \node at (0,-0.3) {\small $\gamma$};
    \node at (4.5,-0.3) {\small $\gamma$};
\end{tikzpicture}
\end{center}
The $\varepsilon^{\mu\nu\rho\sigma}\varepsilon_{\alpha\beta\gamma\delta}$ contraction from two 
$F\tilde{F}$ vertices produces a factor 2 from the antisymmetric tensor algebra:
\begin{equation}
    \varepsilon^{\mu\nu\rho\sigma}\varepsilon_{\mu\nu\alpha\beta} = -2(\delta^\rho_\alpha\delta^\sigma_\beta - \delta^\rho_\beta\delta^\sigma_\alpha)
\end{equation}

\textbf{(ii) Factor 8 and $\cthree^6$ in $8\bone\cthree^6 L$:}

The log-dependent term has two components:

\textbf{Origin of $\cthree^6$:} The factor $\cthree^6$ arises from the product structure 
of the effective action. With $A = 2\cthree^3$ as the overall amplitude:
\begin{equation}
    8\bone\cthree^6 L = 8\bone \cdot \cthree^3 \cdot \cthree^3 \cdot L
    = 4 \cdot A \cdot \cthree^3 \cdot \bone L
\end{equation}
The factorization proceeds as follows:
\begin{itemize}[leftmargin=1.5em, topsep=2pt, itemsep=1pt]
    \item One factor of $\cthree^3$ comes from three $\cthree\, aF\tilde{F}$ vertices 
        in the anomaly chain (triangle diagram).
    \item The second factor of $\cthree^3$ comes from the amplitude normalization $A/2 = \cthree^3$.
    \item The coefficient 8 decomposes as $8 = 4 \times 2$, where 4 is from $|D_2|^2$ 
        (two $\mathbb{Z}_2$ reflections) and 2 is from the double cover sheets.
\end{itemize}
\textbf{Consistency check:} $8\cthree^6 = 8 \times (8\pi)^{-6} = 8/(8\pi)^6$, which matches 
the numerical coefficient in the CFE.

\textbf{Origin of factor 8:} The dihedral symmetry $D_4$ of the 4-point amplitude 
contributes $|D_4| = 8$. This arises from the 4 rotations and 4 reflections that 
leave the fermion loop topology invariant:
\begin{center}
\begin{tikzpicture}[scale=0.8]
    \draw[thick] (0,0) rectangle (2,2);
    \draw[thick, photon] (-0.8,0) -- (0,0);
    \draw[thick, photon] (-0.8,2) -- (0,2);
    \draw[thick, photon] (2,0) -- (2.8,0);
    \draw[thick, photon] (2,2) -- (2.8,2);
    \node at (1,-0.5) {\small fermion loop};
    \node at (1,1) {\small $D_4$};
\end{tikzpicture}
\end{center}

\textbf{The $\bone$ factor:} $\bone = 41/10$ counts the hypercharge-weighted fermion 
contributions summed over the SM spectrum. It enters as $\sum_f Y_f^2$, where the 
sum runs over all SM fermions with their multiplicities.

\textbf{(iii) Factor 2/3 in potential:}

The coefficient $2/3$ in the cubic term of \eqref{eq:potential} ensures that 
differentiation produces the factor 2 in \eqref{eq:stationarity}:
\begin{equation}
    \frac{d}{d\alpha}\left(-\frac{2}{3}A\cthree^3\alpha^3\right) = -2A\cthree^3\alpha^2
\end{equation}
This is not a choice but a consequence of the integral representation of the 
effective potential: the cubic term arises from a triangle subgraph.

\textbf{Why no $\alpha^2$ term?}

The absence of an $\alpha^2$ term in the CFE is not accidental. In background-field gauge, 
quadratic terms are absorbed into the field-strength renormalization. The Ward identity 
ensures that physical observables depend only on odd powers of $\alpha$ (gauge invariance 
requires $F^2 \sim \alpha^2$ to enter as a full factor, not as a correction to $\alpha$).

\begin{tcolorbox}[colback=purple!3, colframe=purple!50!black, 
    title={\textbf{Why Integer Coefficients? (No $\zeta(3)$, $\ln 2$, etc.)}}, fonttitle=\bfseries\small]
\textbf{Potential objection:} In typical QFT calculations (e.g., anomalous magnetic moment $g-2$), 
coefficients involve transcendental numbers like $\zeta(3)$, $\pi^2/12$, $\ln 2$. Why are the 
CFE coefficients (2, 8, 48) simple integers?

\textbf{Resolution:} The CFE is not a perturbative loop expansion but a \emph{topological constraint}.

\begin{enumerate}[leftmargin=1.5em, topsep=2pt, itemsep=1pt]
    \item \textbf{Origin of transcendentals:} In $g-2$, transcendentals arise from 
        \emph{integrating over continuous momenta} in Feynman integrals. The integrals 
        produce polylogarithms, which evaluate to $\zeta$-values.
    
    \item \textbf{Topological discreteness:} In TFPT, the relevant integrals are over 
        \emph{discrete topological data}: Euler characteristics, winding numbers, and 
        symmetry group orders. These are integers by definition.
    
    \item \textbf{Combinatorial origin:}
        \begin{itemize}[leftmargin=1em, topsep=1pt, itemsep=1pt]
            \item Factor 2: From $\varepsilon^{\mu\nu\rho\sigma}$ contraction (antisymmetry)
            \item Factor 8: From $|D_4| = 8$ (dihedral group of the box)
            \item Factor 48: \emph{Mnemonic} $2 \times 24$ (double cover $\times$ permutations); the derivation is the spin-lifted deficit with four cut patches (Appendix~\ref{app:48deriv})
        \end{itemize}
    
    \item \textbf{Where $\pi$ appears:} The factors of $\pi$ in $\cthree = 1/(8\pi)$ and 
        $\phitree = 1/(6\pi)$ come from the \emph{normalization} of circles ($2\pi$ per cycle), 
        not from loop integrals. They are geometric, not perturbative.
\end{enumerate}

\textbf{Diagnostic:} If transcendentals appeared in the CFE, it would indicate non-topological 
dynamics. Their absence is evidence that the fixed-point structure is genuinely topological.
\end{tcolorbox}

\begin{remark}
The consistency between \eqref{eq:potential} and \eqref{eq:CFE-appendix} can be 
verified directly: $\frac{d}{d\alpha}\bigl(-\frac{2}{3}A\cthree^3\alpha^3\bigr) 
= -2A\cthree^3\alpha^2$, matching the required coefficient.
\end{remark}

%=============================================================================
\section{Verification Suite Summary}
%=============================================================================

The TFPT verification suite (\texttt{tfpt-suite/}) turns paper claims into deterministic, assumption-explicit computations. All counts in this section are auto-generated from the suite artifacts (no hardcoded ``39 modules''):
\[
\text{unique modules (registry)}=\TFPTSuiteUniqueModulesTotal,\qquad
\text{module artifacts (engineering+physics)}=\TFPTSuiteModulesTotal,\qquad
\text{checks total}=\TFPTSuiteChecksTotal.
\]
It ships:
\begin{itemize}[leftmargin=1.5em, topsep=2pt, itemsep=1pt]
    \item a \textbf{conventional suite} (registry-driven; core invariants, $\alpha$ closure + policy bridge, RG/matching, flavor pipelines, cosmology/bounce wiring, torsion/birefringence, dashboards),
    \item an \textbf{unconventional suite} of search/audit/design tools (for closing explicit ToE gaps without hiding assumptions),
    \item a \textbf{theoryv3} branch of discrete audits that formalizes the ``kernel + crosslinks'' thesis (holonomy $g=5$, defect partition, constant factory).
\end{itemize}
Each module emits \texttt{results.json} (machine-readable), \texttt{report.txt} (human-readable), and optional plots; PDF reports summarize runs.

\subsection{Known Gaps and Status (Auto-Generated)}
\label{sec:known-gaps}

\begin{tcolorbox}[colback=yellow!3, colframe=yellow!60!black,
title={\textbf{Current WARN/FAIL inventory (from suite manifest)}}, fonttitle=\bfseries\small]
\TFPTSuiteProblemModules
\end{tcolorbox}

% Compact table of problem modules (WARN/FAIL/plot issues)
% Inlined (standalone build): formerly \input{../tfpt-suite/out/suite_manifest_table.tex}
% -----------------------------------------------------------------------------
% AUTO-GENERATED by tfpt_suite.suite_manifest.py — do not edit by hand.
\begin{table}[H]
\centering
\small
\begin{tabular}{@{}llrrrrrr@{}}
\toprule
Module & Mode & Checks & PASS & WARN & FAIL & Plots(exp) & Plots(ok)\\
\midrule
\texttt{alpha\_on\_shell\_bridge} & engineering & 7 & 6 & 1 & 0 & 0 & 0\\
\texttt{matching\_finite\_pieces} & engineering & 5 & 4 & 1 & 0 & 0 & 0\\
\texttt{torsion\_condensate} & engineering & 8 & 5 & 3 & 0 & 0 & 0\\
\texttt{alpha\_on\_shell\_bridge} & physics & 7 & 6 & 1 & 0 & 0 & 0\\
\texttt{matching\_finite\_pieces} & physics & 5 & 4 & 1 & 0 & 0 & 0\\
\texttt{torsion\_condensate} & physics & 8 & 5 & 2 & 1 & 0 & 0\\
\bottomrule
\end{tabular}
\caption{Suite manifest: compact table of WARN/FAIL modules and plot issues (auto-generated).}
\end{table}

